\newpage
\section{第23章：全称类型}

\begin{taplauthorsolutionentrybox}
\phantomsection\label{sol:author:23.4.3}
\textbf{23.4.3 解答}

先定义列表连接，再用它反转列表：
\[
\begin{aligned}
\taplterm{append}={}&\lambda X.\,
\taplterm{fix}\bigl(\lambda app:
\tapltype{List}\,X\to\tapltype{List}\,X\to\tapltype{List}\,X.\,
\lambda l_1:\tapltype{List}\,X.\lambda l_2:\tapltype{List}\,X.\\
&\quad\textbf{if }\taplterm{isnil}[X]l_1\textbf{ then }l_2\\
&\quad\textbf{ else }\taplterm{cons}[X](\taplterm{head}[X]l_1)
  (app(\taplterm{tail}[X]l_1)l_2)\bigr),\\
\taplterm{reverse}={}&\lambda X.\,
\taplterm{fix}\bigl(\lambda rev:\tapltype{List}\,X\to\tapltype{List}\,X.\,
\lambda l:\tapltype{List}\,X.\\
&\quad\textbf{if }\taplterm{isnil}[X]l
\textbf{ then }\taplterm{nil}[X]\\
&\quad\textbf{ else }\taplterm{append}[X]\,
  (rev(\taplterm{tail}[X]l))\,
  (\taplterm{cons}[X](\taplterm{head}[X]l)(\taplterm{nil}[X]))\bigr).
\end{aligned}
\]
\taplsolutionbacklink{ex:23.4.3}{返回练习23.4.3}
\end{taplauthorsolutionentrybox}

\begin{taplauthorsolutionentrybox}
\phantomsection\label{sol:author:23.4.5}
\textbf{23.4.5 解答}
\[
\taplterm{and}=\lambda b:\tapltype{CBool}.\lambda c:\tapltype{CBool}.
\lambda X.\lambda t:X.\lambda f:X.\,b[X](c[X]t\,f)f
\]
\taplsolutionbacklink{ex:23.4.5}{返回练习23.4.5}
\end{taplauthorsolutionentrybox}

\begin{taplauthorsolutionentrybox}
\phantomsection\label{sol:author:23.4.6}
\textbf{23.4.6 解答}
\[
\taplterm{iszero}=\lambda n:\tapltype{CNat}.\,
n[\tapltype{Bool}](\lambda b:\tapltype{Bool}.\taplterm{false})
\taplterm{true}
\]
\taplsolutionbacklink{ex:23.4.6}{返回练习23.4.6}
\end{taplauthorsolutionentrybox}

\begin{taplauthorsolutionentrybox}
\phantomsection\label{sol:author:23.4.8}
\textbf{23.4.8 解答}
\[
\begin{aligned}
\taplterm{pairNat}
 &=\lambda n_1:\tapltype{CNat}.\lambda n_2:\tapltype{CNat}.
   \lambda X.\lambda f:\tapltype{CNat}\to\tapltype{CNat}\to X.\,
   f\,n_1\,n_2,\\
\taplterm{fstNat}
 &=\lambda p:\tapltype{PairNat}.\,
   p[\tapltype{CNat}]
   (\lambda n_1:\tapltype{CNat}.\lambda n_2:\tapltype{CNat}.n_1),\\
\taplterm{sndNat}
 &=\lambda p:\tapltype{PairNat}.\,
   p[\tapltype{CNat}]
   (\lambda n_1:\tapltype{CNat}.\lambda n_2:\tapltype{CNat}.n_2).
\end{aligned}
\]
\taplsolutionbacklink{ex:23.4.8}{返回练习23.4.8}
\end{taplauthorsolutionentrybox}

\begin{taplauthorsolutionentrybox}
\phantomsection\label{sol:author:23.4.9}
\textbf{23.4.9 解答}
\[
\begin{aligned}
\taplterm{zz}&=\taplterm{pairNat}\ \taplterm{c0}\ \taplterm{c0},\\
\taplterm{f}&=\lambda p:\tapltype{PairNat}.\,
 \taplterm{pairNat}\,
 (\taplterm{csucc}(\taplterm{fstNat}\ p))\,
 (\taplterm{fstNat}\ p),\\
\taplterm{pred}&=\lambda m:\tapltype{CNat}.\,
 \taplterm{sndNat}\bigl(m[\tapltype{PairNat}]\taplterm{f}\taplterm{zz}\bigr).
\end{aligned}
\]
迭代 \(i\) 次后得到 \((i,i-1)\)（零处截断），故取第二分量即为前驱。
\taplsolutionbacklink{ex:23.4.9}{返回练习23.4.9}
\end{taplauthorsolutionentrybox}

\begin{taplauthorsolutionentrybox}
\phantomsection\label{sol:author:23.4.10}
\textbf{23.4.10 解答}
\[
\begin{aligned}
\taplterm{vpred}={}&\lambda n:\tapltype{CNat}.\lambda X.\lambda s:X\to X.
\lambda z:X.\\
&\bigl(n[(X\to X)\to X]\,
  (\lambda p:(X\to X)\to X.\lambda q:X\to X.\\
&\hspace{29mm}q(p\,s))\,
  (\lambda x:X\to X.\,z)\bigr)(\lambda x:X.\,x)
\end{aligned}
\]
其类型为 \(\tapltype{CNat}\to\tapltype{CNat}\)。迭代过程中保留一个“尚待
应用的后继”，最后以恒等函数收束，从而丢掉第一次迭代、留下前驱结果。
\taplsolutionbacklink{ex:23.4.10}{返回练习23.4.10}
\end{taplauthorsolutionentrybox}

\begin{taplauthorsolutionentrybox}
\phantomsection\label{sol:author:23.4.11}
\textbf{23.4.11 解答}
\[
\taplterm{head}=
\lambda X.\lambda default:X.\lambda l:\tapltype{List}\,X.\,
l[X](\lambda hd:X.\lambda tl:X.\,hd)\,default
\]
\taplsolutionbacklink{ex:23.4.11}{返回练习23.4.11}
\end{taplauthorsolutionentrybox}

\begin{taplauthorsolutionentrybox}
\phantomsection\label{sol:author:23.4.12}
\textbf{23.4.12 解答}

关键是让折叠同时构造两份列表：第一份尚未插入 \(e\)，第二份已经把 \(e\)
插入适当位置。自右向左处理元素 \(hd\) 时，比较 \(e\) 与 \(hd\)：若
\(e\leq hd\)，把 \(e\) 插在当前后缀之前；否则沿用后缀中已经完成的插入。
\[
\begin{aligned}
\taplterm{insert}={}&\lambda X.\lambda leq:X\to X\to\tapltype{Bool}.
\lambda l:\tapltype{List}\,X.\lambda e:X.\\
&\textbf{let }res=
l[\tapltype{Pair}(\tapltype{List}\,X)(\tapltype{List}\,X)]\\
&\quad(\lambda hd:X.\lambda acc:\tapltype{Pair}
(\tapltype{List}\,X)(\tapltype{List}\,X).\\
&\qquad\textbf{let }rest=\taplterm{fst}\ acc\textbf{ in}\\
&\qquad\textbf{let }rest_e=\taplterm{snd}\ acc\textbf{ in}\\
&\qquad\taplterm{pair}\,
  (\taplterm{cons}[X]hd\,rest)\,
  (\textbf{if }leq\ e\ hd\\
&\hspace{52mm}\textbf{ then }\taplterm{cons}[X]e
     (\taplterm{cons}[X]hd\,rest)\\
&\hspace{52mm}\textbf{ else }\taplterm{cons}[X]hd\,rest_e))\\
&\quad(\taplterm{pair}(\taplterm{nil}[X])
  (\taplterm{cons}[X]e(\taplterm{nil}[X])))\\
&\textbf{ in }\taplterm{snd}\ res.
\end{aligned}
\]
随后再以插入函数折叠整个列表：
\[
\begin{aligned}
\taplterm{sort}={}&\lambda X.\lambda leq:X\to X\to\tapltype{Bool}.\\
&\lambda l:\tapltype{List}\,X.\,
l[\tapltype{List}\,X]\\
&\quad
(\lambda hd:X.\lambda rest:\tapltype{List}\,X.\,
\taplterm{insert}[X]leq\,rest\,hd)(\taplterm{nil}[X]).
\end{aligned}
\]
\taplsolutionbacklink{ex:23.4.12}{返回练习23.4.12}
\end{taplauthorsolutionentrybox}

\begin{taplauthorsolutionentrybox}
\phantomsection\label{sol:author:23.5.1}
\textbf{23.5.1 解答}

证明结构几乎与\cref{thm:9.3.9}相同。E-TAppTabs 情形还需要类型替换引理：
\[
\Gamma,X,\Delta\vdash t:T
\quad\Longrightarrow\quad
\Gamma,[X\mapsto S]\Delta\vdash[X\mapsto S]t:[X\mapsto S]T.
\]
保留尾部上下文 \(\Delta\) 才能得到足够强的归纳假设；否则 T-Abs 情形无法
处理在 \(X\) 之后引入的项变量绑定。对求值推导归纳，新增的 E-TApp 情形直接
使用归纳假设，E-TAppTabs 情形使用上述引理。
\taplsolutionbacklink{thm:23.5.1}{返回定理23.5.1}
\end{taplauthorsolutionentrybox}

\begin{taplauthorsolutionentrybox}
\phantomsection\label{sol:author:23.5.2}
\textbf{23.5.2 解答}

证明与简单类型 lambda 演算的进展性证明相同，只需扩充典范形式引理：
\[
\text{若值 }v\text{ 具有类型 }\forall X.T_{12},
\text{ 则 }v=\lambda X.t_{12}.
\]
在 T-TApp 情形中，先对被实例化项使用归纳假设；若它能求值，应用 E-TApp；
若它是值，典范形式引理说明它必为类型抽象，于是 E-TAppTabs 可用。
\taplsolutionbacklink{thm:23.5.2}{返回定理23.5.2}
\end{taplauthorsolutionentrybox}

\begin{taplauthorsolutionentrybox}
\phantomsection\label{sol:author:23.6.3}
\textbf{23.6.3 解答}

证明分七步整理类型抽象、类型应用与普通项的相对位置。
\begin{enumerate}
  \item 对 \(t\) 归纳并使用赋型反演：若 \(t\) 良类型且擦除为 \(m\)，可
        重排出一个同样擦除为 \(m\) 的良类型暴露项。
  \item 对最外层类型抽象和类型应用的数量归纳，把任意装饰写成
        “若干类型抽象、一个暴露项、若干类型应用”的形式。相邻的类型抽象和
        类型应用可以先作一次类型 beta 归约；保持性保证重排后类型不变。
  \item 若暴露项擦除为应用 \(m\,n\)，反演可知它必为 \(s\,u\)，其中
        \(s:U\to T\)、\(u:U\)，且分别擦除为 \(m,n\)。
  \item 把上述结论用于擦除为 \(x\,x\) 的项，得到同一个变量经过两组类型
        实例化后既充当函数又充当实参的必要类型方程。
  \item 若项擦除为 \(\lambda x.m\)，整理后的类型必由若干全称量词包住一个
        箭头类型。
  \item 定义类型的“最左叶”：变量的最左叶是自身，箭头类型取参数的最左叶，
        全称类型取量词体的最左叶。对类型大小归纳可知，步骤 4 的方程迫使
        箭头参数的最左叶等于某个外层量化变量。
  \item 反设 \(\Omega\) 可赋型。分别把两个擦除为
        \(\lambda x.x\,x\) 的子项整理并反演。函数位置要求某一类型的最左叶
        等于在左侧抽象处绑定的类型变量；实参位置又要求同一最左叶等于在右侧
        抽象处绑定的另一个类型变量。两变量由不同绑定子引入，不可能相同，
        矛盾。
\end{enumerate}
因此 \(\Omega\) 在 System F 中不可赋型。
\taplsolutionbacklink{ex:23.6.3}{返回练习23.6.3}
\end{taplauthorsolutionentrybox}

\begin{taplauthorsolutionentrybox}
\phantomsection\label{sol:author:23.7.1}
\textbf{23.7.1 解答}
\[
\textbf{let }r=\lambda X.\,\taplterm{ref}(\lambda x:X.\,x)
\textbf{ in }
\bigl(r[\tapltype{Nat}]:=(\lambda x:\tapltype{Nat}.\taplterm{succ}\,x);
\,!(r[\tapltype{Bool}])\ \taplterm{true}\bigr)
\]
若错误地把同一个可变单元在不同实例中分别视为
\(\tapltype{Nat}\to\tapltype{Nat}\) 和
\(\tapltype{Bool}\to\tapltype{Bool}\)，写入自然数函数后再按布尔函数读取，
就破坏了类型安全。
\taplsolutionbacklink{ex:23.7.1}{返回练习23.7.1}
\end{taplauthorsolutionentrybox}
